\section{Discussion and Limitations}

The main methodological implication of this paper is that benchmark validity should be treated as part of intelligent-vehicle evaluation whenever the driving policy depends on local LLM inference. In a conventional simulation study, ranking is often assumed to be the natural endpoint of the experiment. Our results suggest a stricter interpretation for local language-model agents: before comparing policies, the evaluation must first establish that the runtime regime leaves those policies jointly interpretable. This does not replace scenario coverage, safety validation, or broader vehicle-level testing. It adds an additional layer of measurement discipline for language-driven controllers whose inference behavior can dominate the usefulness of the benchmark itself.

That implication also changes how fine-tuning results should be read in automated-driving research. A fine-tuned policy that improves score but does so only in the region where the evaluation remains stable is scientifically useful; a fine-tuned policy that cannot remain rankable under the intended local runtime regime is harder to justify as an evaluation win. The paper therefore argues for a narrower but more defensible reading of local adaptation: fine-tuning evidence should be attached not only to task improvement, but also to the validity conditions under which that improvement was observed.

The main limitation of the current paper is sparse exact pair coverage. Although the repository supports many model families and multiple tiers, only one exact base-vs-fine-tuned pair is ranking-eligible in the current evidence bundle. This means the fine-tuning story is real but still narrow. The paper should be explicit that family-level generalization remains unresolved.

The second limitation is scope. The repository contains cross-scenario infrastructure for merge and intersection tasks, but those settings are not part of the empirical evidence in the present manuscript. The first paper should therefore remain highway-only and avoid implying any cross-scenario generalization. For intelligent-vehicle evaluation, this also means the current study should be read as a targeted examination of local decision-making behavior in one highway benchmark family rather than as a complete validation protocol for automated driving.

A third limitation is runtime-regime specificity. The current timeout ladder, hardware context, and local serving stack are part of the studied setup. That is a strength for local reproducibility, but it also means the present manuscript should not present the resulting validity boundary as a universal benchmark standard. The paper can claim that runtime validity is a first-class measurement concern for local LLM driving agents; it cannot claim that the exact scale boundary observed here will transfer unchanged to every inference stack or vehicle-evaluation workflow.
